\documentclass[12pt]{article}
\usepackage{mathtools}
\usepackage{amssymb}
\usepackage{fontspec}
% \usepackage[letf=1.06cm,top=0.9cm,right=1.06cm,bottom=0.49cm]{geometry}
% \setmainfont{GFS Didot}
\setmainfont{EB Garamond}


% Set page size and margins
\usepackage[a4paper,top=1.5cm,bottom=1.5cm,left=1.3cm,right=1.3cm]{geometry}
\usepackage{fancyhdr}
\pagestyle{fancy}
\setlength{\parindent}{0pt}

% Useful packages
\usepackage{minted}
\usepackage{graphicx}
\graphicspath{{./media/}}
\usepackage{subfig}
\usepackage{float}
\usepackage[colorlinks=true, allcolors=blue]{hyperref}
\usepackage{siunitx}
\usepackage{bm}
% \usepackage{sectsty}
% \sectionfont{\fontsize{12}{15}\selectfont}

\title{\vspace{-2cm}052 - Τεχνικές Βελτιστοποίησης\\
    \large Project}

\author{Ιωάννης Δημουλιός \\ 10641}
\date{Χειμερινό εξάμηνο 2025}

\lhead{052 - Τεχνικές Βελτιστοποίησης}
\chead{Project}
\rhead{Ιωάννης Δημουλιός}
\begin{document}
\maketitle
\textbf{Εισαγωγή} \\
Θέλουμε να προσεγγίσουμε μια άγωστη, συνεχή συνάρτηση \(y = f(u_1, u_2)\) ως το βεβαρημένο άθροισμα το πολύ $ M \leq 15$ γκαουσσιανών συναρτήσεων 
\[G_k(u_1, u_2) = \exp\left(-\left( \dfrac{(u_1 - c_{1k})^2}{2\sigma_{1k}^2} + \dfrac{(u_2 - c_{2k})^2}{2\sigma_{2k}^2}\right) \right)\] 
με κέντρα και τυπικές απόκλισεις ανά άξονα $\mu_{k1}$, $\mu_{k2}$ και $\sigma_{k1}$, $\sigma_{k2} > 0$ αντίστοιχα, όπου $1 \leq k \leq M$. Επομένως, ψάχνουμε μια συνάρτηση 
\[
\hat{f}(u_1, u_2) = \sum_{k=1}^{M} w_k G_k(u_1, u_2)
.\] 
που να ελαχιστοποιεί την απόσταση (που θα ορίσουμε παρακάτω) $d(f, \hat{f})$. Δηλαδή το πρόβλημα βελτιστοποίσης γράφεται 
\[
\min_{\substack{w_k, \mu_{k1}, \mu_{k2} ,\sigma_{k1}, \sigma_{k2}\\ 1 \leq k \leq M}} d(f, \hat{f})
\]
και αναζητούμε τα κατάλληλα βάρη, κέντρα και τυπικές αποκλίσεις. 
Για να το πετύχουμε αυτό έχουμε ένα πεπερασμένο σύνολο εξάσκησης $N$ σημείων της συνάρτησης 
\[
\{(u_{1j}, u_{2j}, y_j),\: y_j = f(u_{1j}, u_{2j}), \: 1 \leq j \leq N\}.
\] \\

\textbf{Συνάρτηση Απόστασης} \\
Αν \(\hat{y}_j = \hat{f}(u_{1j}, u_{2j})\), τότε ορίζουμε τη συνάρτηση απόστασης (σφάλματος) μεταξύ της αρχικής και της εκτιμήτριας συνάρτησης ως το άθροισμα των τετραγωνικών σφαλμάτων
\[
d(f, \hat{f}) = \sum_{j=1}^{N} (y_j - \hat{y}_j)^2
.\]

\textbf{Γενετικός Αλγόριθμος} \\
Για την υλοποίηση της βελτιστοποίησης χρησιμοποιούμε γενετικό Αλγόριθμο με τα εξής ορίσματα και εξόδους
\begin{verbatim}
[best_individual, best_error] = genetic_approximator([u1j], [u2j], [yj], M)
\end{verbatim}
Ένας \texttt{individual} είναι ένα διάνυσμα $5M$ στοιχείων που περιέχει τις παραμέτρους $(w_k, c_{1k}, c_{2k}, \sigma_{1k}, \sigma_{2k})$ για τις $M$ γκαουσσιανές $G_k$.\\
Ο \verb|best_individual| είναι ένα διάνυσμα $5M$ στοιχείων που περιέχει τις παραμέτρους $(w_k^*, c_{1k}^*, c_{2k}^*, \sigma_{1k}^*, \sigma_{2k}^*)$ για τις $M$ γκαουσσιανές $G_k^*$ που παράγουν τη βέλτιστη εκτιμήτρια $\hat{f}^*$ και ο \verb|best_error| είναι το βέλτιστο σφάλμα $d(f, \hat{f}^*)$. \\

\textit{Αρχικοποίηση Παραμέτρων Βελτιστοποίσης} \\
Αρχικοποιούμε έναν πληθυσμό από \texttt{individual} χρησιμοποιώντας ομοιόμορφες κατανομές σε κατάλληλα διάστηματα που εξαρτώνται από τα εύρη των τιμών των $\{u_{1j}\}$, $\{u_{2j}\}$, $\{y_{j}\}$, 
\[
\ell_1 = \max(u_{1j}) - \min(u_{1j}),\quad \ell_2 = \max(u_{2j}) - \min(u_{2j}), \quad \ell_y = \max(y_j) - \min(y_j)
.\] 
\newpage
\textit{Αξιολόγηση Πληθυσμού} \\
Αν $f'$ η εκτιμήτρια που προκύπτει από έναν συγκεκριμένο \texttt{individual}, τότε χρησιμοποιούμε ως συνάρτηση αξιολόγησης για την καταλληλότητα αυτού του \texttt{individual} την 
\[
v(f') = \dfrac{1}{1 + d(f, f')}
.\]
Τότε έχουμε ότι $v(f') \in (0, 1)$ και όσο καταλληλότερη είναι η $f'$, δηλαδή όσο μικρότερη είναι η $d(f, f')$, τόσο υψηλότερη είναι η $v(f')$. \\

\textit{Επιλογή Γονέων} \\
Χρησιμοποιούμε τη μέθοδο του τροχού της τύχης, όπου η πιθανότητα ο κάθε \texttt{individual} να είναι γονιός είναι 
\[
p(f') = \dfrac{v(f')}{\sum v}
.\] \\ 

\textit{Διασταύρωσεις} \\
Χρησιμοποιούμε διασταύρωση τύπου crossover, δηλαδή επιλέγουμε δύο γονείς \texttt{individual} βάσει της παραπάνω κατανομής και έπειτα επιλέγουμε τυχαία έναν δείκτη όπου κόβουμε τον έναν \texttt{individual} σε δύο κομμάτια και αντικαθιστούμε το ένα κομμάτι με το αντίστοιχο του άλλου \texttt{individual}. Με αυτόν τον τρόπο μπορούμε να φτιάξουμε από δύο διαφορετικούς γονείς δύο διαφορετικά παιδιά. \\
Ο δείκτης όπου πραγματοποιείται η τομή δεν μπορεί να είναι εντελώς τυχαίος· φροντίζουμε να μην κόβει καμιά από τις πεντάδες 
$(w_k, c_{1k}, c_{2k}, \sigma_{1k}, \sigma_{2k})$
ενός \texttt{individual} στη μέση, ώστε να μη μπλέκονται τα μοναδικά χαρακτηριστικά που έχει αναπτύξει η εκάστοτε γκαουσσιανή $G_k$ του \texttt{individual}. 
Η επιλογή του δείκτη τομής γίνεται με τυχαία επιλογή εναποθέτοντας ίσες πιθανότητες σε όλους τους δείκτες που πληρούν το παραπάνω κριτήριο. \\

\textit{Μεταλλάξεις}\\
Κάθε παράμετρος από τις $5M$ ενός παιδιού \texttt{individual} μεταλλάσσεται με σταθερή πιθανότητα. Στις μεταλλασσόμενες παραμέτρους προσθέτουμε μια ποσότητα που παράγεται από γκαουσσιανή κατανομή πολλαπλασιασμένη κατά σταθερό μέγιστο πλάτος. Επομένως, αν \(x\) μια δεδομένη παράμετρος, τότε 
\[
x' = x + bA \Delta x
,\]
όπου $b \in \{0, 1\}$, $Α$ το σταθερό πλάτος και $\Delta x$ ποσότητα που προκύπτει από πείραμα γκαουσσιανής κατανομής. \\
Δοκιμάσαμε και άλλες τεχνικές, όπως για παράδειγμα να μειώνουμε προοδευτικά το πλάτος όσο προχωρούν οι γενιές, αλλά δεν έδιναν καλύτερα αποτελέσματα, οπότε κρατήσαμε την απλούστερη εκδοχή. \\

\textit{Ελιτισμός} \\
Σε κάθε γενιά κρατάμε τον καταλληλότερο \texttt{individual} και αντικάθιστούμε με αυτόν τον χειρότερο παιδί \texttt{individual} στην επόμενη γενιά. Αυτή η τεχνική έδινε καλύτερα αποτελέσματα. \\

\textbf{Παραγωγή Συνόλων Σημείων Εξάσκησης και Αξιολόγησης} \\
Για να ελέγξουμε την ορθότητα του γενετικού αλγορίθμου χρησιμοποιούμε τη συνάρτηση 
\[
f(u_1, u_2) = \sin(u_1 + u_2) \sin(u_2^2)
,\]
με $u_1 \in [-1, 2]$ και $u_2 \in [-2, 1]$. \\
Παράγουμε πρώτα ένα σύνολο εξάσκησης λαμβάνοντας τις τιμές των $u_1$ και $u_2$ τυχαία με ομοιόμορφη κατανομή στο πεδίο ορισμού τους. Το σύνολο αυτό χρησιμοποιείται ως είσοδος στον γενετικό αλγόριθμο
\[
\{(u_{1j}, u_{2j}, y_j),\: y_j = f(u_{1j}, u_{2j}), \: 1 \leq j \leq N\}.
\]
Έπειτα έχοντας τη βελτιστοποιημένη προσέγγιση $f^*$ για το παραπάνω σύνολο εισόδου, την αξιολογούμε παράγοντας ένα νέο σύνολο αξιολόγησης, το οποίο επίσης παράγουμε τυχαία με τον ίδιο τρόπο, 
\[
\{(u'_{1j}, u'_{2j}, y'_j),\: y_j' = f(u'_{1j}, u'_{2j}), \: 1 \leq j \leq N\}.
\]
Έπειτα υπολογίζουμε το νέο σφάλμα
\[
d'(f, f^*) = \sum_{j=1}^{N} (y_j' - y'^*_j)^2
,\]
όπου \(y'^*_j = f^*(u'_{1j}, u'_{2j})\). 

Αν τα $d(f, f^*)$ και $d'(f, f^*)$ είναι συγκρίσιμα και επαρκώς μικρά, τότε θεωρούμε ότι ο αλγόριθμος λειτουργεί σωστά και δίνει καλά αποτελέσματα στο πεδίο ορισμού της $f$. \\

\textbf{Αποτελέσματα} \\
Τρέχουμε τον αλγόριθμο προσπαθώντας να προσεγγίσουμε τη δοθείσα $f$, ορισμένη στο αρχείο \texttt{func.m}, με τις εξής παρακάτω παραμέτρους. Οι πρώτες δύο ορίζονται στο αρχείο \texttt{main.m} και οι υπόλοιπες στο \verb|genetic_approximator.m|:
\begin{itemize}
    \item Πληθάριθμος συνόλου εξάσκησης και αξιολόγησης: $500$, 
    \item Πλήθος γκαουσσιανών: $M=15$,
    \item Μέγεθος πληθυσμούς: $200$, 
    \item Μέγιστος αριθμός γενεών: $12000$,
    \item Πιθανότητα μετάλλαξης παραμέτρου: $0.05$,
    \item Μέγιστο πλάτος μετάλλαξης: $Α=0.1$.
\end{itemize}
    % POPULATION_SIZE = 200; 
    % MAX_GENERATIONS = 12000; 
    % MUTATION_RATE = 0.05; 
    % MUTATION_STRENGTH = 0.1; 
    % RELATIVE_WEIGHT_BOUNDS = [-4 4]; 
    % RELATIVE_STDV_BOUNDS = [0.03 0.3]; 
    % CENTER_TOL = 0.05;
    % TARGET_ERROR = 0.1; 
Λαμβάνουμε, λοιπόν, τη βελτιστοποιημένη προσέγγιση $f^*$. Για αυτήν ισχύει $d(f, f^*) = 0.69$ (αθροιστικό τετραγωνικό σφάλμα ως προς το σύνολο εξάσκησης) από όπου προκύπτει ρίζα μέσου τετραγωνικού σφάλματος (rmse) 
$$e = \sqrt{d(f, f^*)/N} = 0.037.$$ 
Ακόμη $d'(f, f^*) = 0.91$ (αθροιστικό τετραγωνικό σφάλμα ως προς το σύνολο αξιολόγησης). \\
Έτσι, παίρνουμε rmse $e' = 0.042$. 
Τα δεδομένα αυτά θεωρούνται ικανοποιητικά, αφού το σύνολο τιμών της συνάρτησης είναι το $[-1, 1]$. \\

\textbf{Παράρτημα} \\
\textit{Οδηγίες χρήσης του κώδικα} \\
Το σύνολο του κώδικα βρίσκεται μέσα στον φάκελο \texttt{src/}. \\
Αν θέλουμε να προσεγγίσουμε μια γνωστή συνάρτηση, όπως στην άσκηση, την ορίζουμε στο \texttt{func.m}. \\
Οι παράμετροι λειτουργίας του γενετικού αλγορίθμου μπορούν να τροποποιηθούν στο \verb|genetic_approximator.m|. Μόνο το πλήθος των χρησιμοποιούμενων γκαουσσιανών ορίζεται στο \texttt{main.m} ως \verb|GAUSSIAN_NUM|.  \\
Το \verb|dataset_generator.m| περιέχει τη συνάρτηση για την παραγωγή των συνόλων εξάσκησης και αξιολόγησης. \\
Το \verb|reconstructor.m| περιέχει τη συνάρτηση που δοθέντος ενός \texttt{individual} και ενός συνόλου τιμών εισόδου επιστρέφει τις τιμές εξόδου της προσέγγισης που αντιστοιχεί στον \texttt{individual}. \\
Το \verb|evaluator.m| επιστρέφει το αθροιστικό τετραγωνικό σφάλμα μιας προσέγγισης που αντιστοιχεί σε \texttt{individual} ως προς ένα σύνολο πραγματικών τιμών εξόδου της συνάρτησης (π.χ. σύνολο εξάσκησης και αξιολόγησης). 
\newpage
\textit{Διαγράμματα σε όλο το πεδίο ορισμού της συνάρτησης}
\begin{figure}[H]
    \centering
    \includegraphics[width=0.9\textwidth]{3d_comparison.png}
    Διάγραμμα αρχικής συνάρτησης (αριστερά) και προσέγγισης (δεξιά) \\
    Ο γενετικός αλγόριθμος έχει βρει και έχει προσεγγίσει ικανοποιητικά όλους τους λόφους και τις κοιλάδες της συνάρτησης.
\end{figure}
\begin{figure}[H]
    \centering
    \includegraphics[width=0.9\textwidth]{error_heatmap.png}
    Διάγραμμα απόλυτου σφάλματος εξόδου προσέγγισης ως προς αρχική συνάρτηση\\
    Το σφάλμα είναι ικανοποιητικά μικρό στο μεγαλύτερο μέρος του πεδίου ορισμού. Μόνο στα άκρα του σε κάποια σημεία είναι εμφανώς μεγαλύτερο. Αυτό οφείλεται στην έμφυτη τυχαιότητα κατά τη δημιουργία του συνόλου εξάσκησης το οποίο μάλλον δεν κάλυπτε τα άκρα του πεδίου οριμού στο σύνολό τους. 
\end{figure}

\end{document}